\subsection{Formulação dinâmica acoplada}
\label{sec:dinamica_acoplada}

A dinâmica do sistema base--manipulador considera que o conjunto opera como um sistema mecânico isolado, sujeito às forças e torques internos, sendo assim, o comportamento dinâmico é obtido a partir da energia cinética total, composta pelas contribuições da base e dos elos do MR:

\begin{equation}
\label{eq:T_total}
T \;=\; T_{B} \;+\; \sum_{i=1}^{5} T_{i},
\end{equation}

onde $T_B$ é a energia cinética translacional e rotacional da base e $T_i$ é a energia cinética do $i$-ésimo elo, expressas no referencial da base $\{B\}$.

A energia cinética da base é dada por:

\begin{equation}
\label{eq:TB}
T_{B} \;=\;
\frac{1}{2} \, ({}^{B} \vec v_{B}^{\,T}) (m_{B}) I_{3} \, ({}^{B}\,\vec v_{B})
\;+\;
\frac{1}{2} \, ({}^{B}  \, \vec \omega_{B}^{\,T}) \, I_{B} \, ({}^{B}\,\vec \omega_{B}),
\end{equation}

onde $m_B$ é a massa da base e $I_B$ é o tensor de inércia da base em torno do seu centro de massa.
Para cada elo do manipulador, a energia cinética pode ser expressa por: 

\begin{equation}
\label{eq:Ti}
T_i \;=\;
\frac{1}{2} \, m_i \, ({}^{B} \,\vec v_{G_i}^{\,T}) \, {}^{B} \,\vec v_{G_i}
\;+\;
\frac{1}{2} \, ({}^{B} \,\vec \omega_i^{\,T}) \, I_i \, ({}^{B} \,\vec \omega_i),
\end{equation}

onde ${}^{B} \,\vec v_{G_i}$ e ${}^{B} \,\vec \omega_i$ são, respectivamente, a velocidade linear do centro de massa e a velocidade angular do elo $i$, ambos expressos no referencial da base $\{B\}$.
As velocidades ${}^{B} \,\vec v_{G_i}$ e ${}^{B} \,\vec \omega_i$ dependem das velocidades generalizadas da base e das juntas e são obtidas a partir do Jacobiano acoplado da Seção~\ref{sec:cinematica_velocidades}.
Substituindo \eqref{eq:twist_coupled} em \eqref{eq:T_total}, a energia cinética total pode ser escrita como:

\begin{equation}
\label{eq:T_matriz}
T \;=\; \frac{1}{2} \, \dot{\vec q}^{\,T} \, M(\vec q) \, \dot{\vec q},
\qquad
\dot{\vec q}^{\,T} \;=\; 
\big[\, \dot{\vec q}_{B}^{\,T} \; \dot{\vec\theta}^{\,T} \,\big],
\end{equation}

onde $M(\vec q)$ é a matriz de inércia generalizada do sistema.

A derivação do Lagrangiano e a aplicação das equações de Euler--Lagrange conduzem à equação dinâmica geral:

\begin{equation}
\label{eq:dinamica_geral}
M(\vec q)\,\ddot{\vec q}
\;+\;
C(\vec q, \dot{\vec q})\,\dot{\vec q}
\;=\;
\vec\tau
\;+\;
\vec\tau_{\text{int}},
\end{equation}

onde $C(\vec q, \dot{\vec q})$ contém os termos centrífugos e de Coriolis, $\vec\tau$ são os torques e forças aplicadas pelas juntas, e $\vec\tau_{\text{int}}$ contém os torques de acoplamento interno entre base e manipulador.
A matriz de inércia $M(\vec q)$ possui a seguinte estrutura:

\begin{equation}
\label{eq:M_blocos}
M(\vec q) \;=\;
\begin{bmatrix}
M_{bb} & M_{bm} \\[4pt]
M_{mb} & M_{mm}
\end{bmatrix},
\end{equation}

em que:
\begin{itemize}
    \item $M_{bb}$ é o bloco associado às inércias translacional e rotacional da base;
    \item $M_{mm}$ é o bloco correspondente às inércias das juntas e elos do manipulador;
    \item $M_{bm}$ e $M_{mb}=M_{bm}^T$ representam o acoplamento dinâmico base--manipulador, e reflete o efeito das movimentações das juntas sobre a atitude da base e vice-versa.
\end{itemize}

O acoplamento pode ser obtido ao separar a equação \eqref{eq:dinamica_geral} em dois conjuntos:

\begin{align}
\label{eq:dinamica_blocos}
M_{bb}\,\ddot{\vec q}_B \;+\; M_{bm}\,\ddot{\vec\theta}
\;+\; C_{bb}\,\dot{\vec q}_B \;+\; C_{bm}\,\dot{\vec\theta}
\;&=\; \vec\tau_B, \\[4pt]
M_{mb}\,\ddot{\vec q}_B \;+\; M_{mm}\,\ddot{\vec\theta}
\;+\; C_{mb}\,\dot{\vec q}_B \;+\; C_{mm}\,\dot{\vec\theta}
\;&=\; \vec\tau_m,
\end{align}

onde $\vec\tau_B$ e $\vec\tau_m$ são os torques e forças aplicadas à base e às juntas, respectivamente.
Por estar voando em um ambiente de microgravidade, não há torque externo resultante sobre o sistema, portanto, o momento angular total é conservado:

\begin{equation}
\label{eq:momentum_total}
\dot{\vec H}_{\text{total}} \;=\; 0,
\qquad
\vec H_{\text{total}} \;=\; M_{bb}\,\dot{\vec q}_B \;+\; M_{bm}\,\dot{\vec\theta}.
\end{equation}

Sendo assim, o movimento das juntas gera reações internas que alteram a atitude da base, sendo assim, mesmo que não tenha propulsão, o controle das juntas pode ser empregado para manobras finas de atitude do conjunto base--manipulador.